%% LyX 2.4.4 created this file.  For more info, see https://www.lyx.org/.
%% Do not edit unless you really know what you are doing.
\documentclass[english]{article}
\usepackage[T1]{fontenc}
\usepackage[latin9]{inputenc}
\usepackage{enumitem}
\usepackage{amsmath}
\usepackage{amsthm}
\usepackage{geometry}
\geometry{verbose,tmargin=1in,bmargin=1in,lmargin=1in,rmargin=1in}

\makeatletter
%%%%%%%%%%%%%%%%%%%%%%%%%%%%%% Textclass specific LaTeX commands.
\numberwithin{equation}{section}
\numberwithin{figure}{section}
\newlength{\lyxlabelwidth}      % auxiliary length 

%%%%%%%%%%%%%%%%%%%%%%%%%%%%%% User specified LaTeX commands.
\usepackage{fancyhdr}
\usepackage{lastpage}
\pagestyle{fancy}
\usepackage{enumitem}
\usepackage{ifthen}
%sets math fonts to be more readable
\everymath=\expandafter{\the\everymath\displaystyle}
%\DeclareMathSizes{10}{10}{10}{10}
\usepackage{multicol}

\usepackage{tikz}
\usetikzlibrary{patterns}
\usetikzlibrary{plotmarks}
\usepackage{pgfplots}
\pgfplotsset{compat=newest}

\usepackage{advdate}    % Advancing/saving dates
\usepackage{datetime}   % Dates formatting
\usepackage{datenumber} % Counters for dates
\usdate
% Added by lyx2lyx
\setlength{\parskip}{\smallskipamount}
\setlength{\parindent}{0pt}

\makeatother

\usepackage{babel}
\begin{document}
\renewcommand{\author}{Dr.\,James Ashe}

\newcommand{\classNum}{232}

\newcommand{\classSec}{A and B}

\newcommand{\nameType}{}

\newcommand{\examType}{Exam 2}

\renewcommand{\date}{will be be given on \formatdate{11}{3}{2013}} %%\currenttime, \today

%%First page's header%%%%%%%%%%%%%%%%%%%%%%
\fancypagestyle{firststyle} %%header settings for first pagr
{
	\fancyhead[L]{MTH \classNum{}(\classSec{}) \author{}\\
	\textbf{\examType{}} \date}
	\fancyhead[C]{\textbf{\nameType}}
	\setlength{\headsep}{14pt}
}
%%%%%%%%%%%%%%%%%%%%%%%%%%%%%%%%%%%%%%%%%%

%%Next page's header%%%%%%%%%%%%%%%%%%%%%%%
\setlist{nolistsep}
\lhead{\classNum{}(\classSec{})} 
\chead{\textbf{\nameType}} 
\cfoot{\thepage{}~of~\pageref{LastPage}} 
\setlength{\headsep}{5pt} 
%%%%%%%%%%%%%%%%%%%%%%%%%%%%%%%%%%%%%%%%%%%

%%Various settings; see the preamble for other settings%%%%%
%%\setenumerate[0]{leftmargin=12pt,itemindent=0pt,labelwidth=0pt}
\renewcommand{\labelenumi}{\textbf{\arabic{enumi}.}}
\renewcommand{\labelenumii}{\textbf{\alph{enumii}.}}
\thispagestyle{firststyle}
%%%%%%%%%%%%%%%%%%%%%%%%%%%%%%%%%%%%%%%%%%

\vfill$ $

\subsection*{5.5 Substitution}
\begin{itemize}
\item Calculating definite integrals using substitution. \textbf{HW 1-3,
6-7}
\item This concept may be partially tested in problems from 6.1-6.3 and
6.5.
\item Problems requiring substitution will not be label ``use substitution''.
Recognizing when to employ this technique is part of the problem. 
\end{itemize}

\subsection*{6.1: Applications of Integration}
\begin{itemize}
\item Given a derivative function, find the displacement or distance traveled
over a given interval. \textbf{HW 1-2, 6, 9-10}
\item Given a derivative function and initial position, find the displacement/position
function. \textbf{HW 3-5, 8}
\end{itemize}

\subsection*{6.2: Area bounded by curves}
\begin{itemize}
\item Find the area between curves. \textbf{HW 1-10}
\item This may require finding the intersection points and using a sum of
two or more integrals.
\end{itemize}

\subsection*{6.3: Volume of solids}
\begin{itemize}
\item General slicing method. \textbf{HW 1-2}

\begin{itemize}
\item Understand the described shape.
\item Place the object in the Cartesian plane such that the cross sections
are lines perpendicular to the \emph{x}-axis (or \emph{y}-axis), and
the height of these lines is determined by a function.
\item Determine an area function for the cross sections dependent on \emph{x
}(or \emph{y}). 
\item Write a definite integral using this area function.
\end{itemize}
\item Disk/Washer method. \textbf{HW 2-8}
\item Applying the FTC, and understanding the inverse relationship between
differentiation and integration. \textbf{HW 1-5, 17-18}
\end{itemize}

\subsection*{6.5 Length of a curve}
\begin{itemize}
\item Use the formula $L={\displaystyle \int_{a}^{b}\sqrt{1+f'(x)^{2}}\,dx}$
to find the length of $f(x)$ from $[a,b]$. \textbf{HW 1-7}
\end{itemize}

\subsection*{Some general notes.}
\begin{itemize}
\item You will be given the necessary formulas from section 4.8. 
\item So that all the concepts can be covered in a 50 minute exam, problems
will generally be computationally easier than the homework OR you
will only be asked to find the definite integral giving the area/volume. 
\item The quiz and homework are good reviews. You can review pass due HW
assignments in MyMathLab (MML) by going to ``Gradebook'' and clicking
on ``Review''. Here you can practice more problems by clicking on
``Similar Exercise''.\vfill
\end{itemize}

\end{document}
